\begin{table*}[!htbp]
\centering
\caption{Costul software al calibrarii SIR (Romania, valul 1; fereastra de $154$ zile, restart-uri multiple Nelder--Mead pentru scenariile liber si constrans). Coloana \emph{in-process} masoara direct costul procedurii din \texttt{calibrate.py}; coloana \emph{naiv (ca job-uri)} proiecteaza costul daca fiecare simulare candidata ar fi un job separat prin aplicatie, folosind latenta end-to-end masurata a unui job SIR real ($1047$~ms/job); coloana \emph{recomandat} corespunde calibrarii intr-un singur job, cu o singura simulare finala scrisa in baza de date. Numarul de simulari candidate este numarul real de evaluari ale functiei obiectiv.}
\label{tab:calibration_cost}
\small
\begin{tabularx}{\textwidth}{X r r r}
\toprule
\textbf{Marime} & \textbf{In-process} & \textbf{Naiv (ca job-uri)} & \textbf{Recomandat} \\
\midrule
Simulari candidate & 1377 & 1377 & 1377 \\
Job-uri prin aplicatie & 1 & 1377 & 1 \\
Randuri scrise in DB & 0 & 213435 & 155 \\
\midrule
Timp total wall (s) & 17.24 & 1441.8\textsuperscript{*} & 17.24 \\
Timp CPU (s) & 17.23 & --- & 17.23 \\
Memorie de varf (MB) & 0.7 & --- & 0.7 \\
\bottomrule
\end{tabularx}
\\[2pt]\footnotesize\textsuperscript{*}~Proiectat: $1377 \times 1047$~ms${}=1442$~s; de $\sim$84$\times$ mai scump decat calibrarea in-process, fara castig de acuratete --- overhead-ul de coada si persistare ar fi platit pentru fiecare integrare candidata.
\end{table*}
